\Qhead{3.6}{Марковские и диффузионные процессы. Броуновское движение. Флуктуационно-диссипационная теорема}{Марковские и диффузионные процессы. Уравнение Фоккера-Планка. Броуновское движение. Флуктуационно-диссипационная теорема.}
\Subfile{3.6.1}{sections/03_statistics/q06/s01}
\Subfile{3.6.2}{sections/03_statistics/q06/s02}

\recap{Марковский процесс описывается текущим состоянием и переходной плотностью. Диффузионный предел даёт Фоккер–Планк. В равновесии ФДТ связывает спектр флуктуаций с диссипативной частью отклика.}
\board{Уравнение Ито, вывод Фоккера–Планка через пробную функцию; броуновский осциллятор; классическая ФДТ.}
\oral{Почему затухание броуновской скорости не приводит к полному исчезновению её равновесных флуктуаций?}
\bait{Можно оставить условие равновесия: произвольный шум активной системы не обязан определяться её потерями по формуле ФДТ.}
\sources{Сёрккя и Солин~\cite{sarkka}, §§5.2--5.4 и глава~6; Clerk и соавт.~\cite{clerk}, §II, формула~(2.16). Пример броуновского движения дополнительно выведен в согласованных единицах.}
